\chapter{Keeping Your Work}
\label{ch:keeping}

\standfirst{A running Smalltalk is a heap of live objects. Three different
things persist it, they persist different things, and knowing which is which
is the difference between losing an afternoon and not.}

\section{The three}

\begin{center}
\small
\begin{tabular}{@{}lp{0.60\textwidth}@{}}
\toprule
\textbf{the image} & the whole world: every object, every window's position,
  your workspace text, the state of everything running. Binary\\[3pt]
\textbf{the changes file} & a text log of every method you accept and every
  expression you \ctp{do it}, appended as you work\\[3pt]
\textbf{file-outs} & source for one class or category, in chunk format, which
  you can read, keep in version control, and file into another image\\
\bottomrule
\end{tabular}
\end{center}

The image is the one that keeps everything and the one you cannot diff.
File-outs are the one you can. Packages on disk---Chapter~\ref{ch:packages}
---are the modern answer and the one to prefer for anything you want to
keep.

\section{Saving the image}

\ct{save} on the desktop menu prompts for a name (defaulting to
\ct{snapshot}), writes \ct{name.image} and \ct{name.changes}, and keeps
running. \ct{quit} offers to do the same and then exit.

From code:

\begin{st}
Smalltalk saveAs: 'myimage' thenQuit: false
\end{st}

Resume it later with:

\begin{sh}
./st80 -run myimage.image
\end{sh}

Everything comes back: windows where you left them, text you had not accepted,
processes that were running.

\begin{stnote}[One arithmetic quirk after a save]
Saving shrinks the display to a hundred rows so the snapshot is small, then
puts the height back. \ct{DisplayScreen} rounds a Form's width down to a
sixteen-bit word boundary on the way, so a 1000-pixel-wide screen comes back
992 wide and the window letterboxes the difference. Harmless, and it is the
image's own arithmetic rather than the VM's.
\end{stnote}

\section{The changes file}

Every method you accept, every class you define and every expression you
\ct{do it} is appended, as text, to \ct{<image>.changes} before it takes
effect. That has two consequences worth knowing.

\textbf{It is your undo of last resort.} If the image is lost, the changes
file still has the source of everything you did, in order, in a format you can
read in any editor and file back in.

\textbf{It is where method source comes from.} The Browser does not store
source in the image---it stores a pointer into a source file, and reads the
text back when you select a method. That is why the changes file must travel
with the image, and why it is created beside it automatically.

\begin{stwarn}
Move an image without its changes file and the Browser will show you nothing
for every method you wrote. The methods still run; their \emph{source} was in
the file you left behind.
\end{stwarn}

\section{Filing out}

\ct{file out} appears on the category, class and message menus in the Browser.
It writes a \ct{.st} file in chunk format:

\begin{center}
\small
\begin{tabular}{@{}ll@{}}
\toprule
category menu & \ctp{Category-Name.st}, every class in it\\
class menu & \ctp{ClassName.st}, that class, both sides\\
message menu & the one method\\
\bottomrule
\end{tabular}
\end{center}

Read one back with \ct{file in} in the File List, or from code:

\begin{st}
(FileStream fileNamed: 'Account.st') fileIn
\end{st}

Chunk format is 1983's interchange format: expressions separated by
exclamation marks, with method bodies in bang-terminated chunks. It is what
everything in \ct{sources/} is written in. It works, and it is worth knowing
because it is how you rescue code from an image.

\section{What to do instead}

For anything you intend to keep, write it as a package on disk---a directory
of Tonel files with a \ct{package.st}---and name it in a profile. Then your
code is:

\begin{itemize}
\item plain text, one file per class, that \ct{git diff} can show you;
\item loadable in any order, because definitions are read in a pass of their
  own;
\item testable from the command line without a window;
\item reproducible, because the profile says exactly what composes the image.
\end{itemize}

The workflow that follows from this is: explore in the image, and when
something works, move it to a file. The Browser is the fastest place to try an
idea and the worst place to leave one.

Chapter~\ref{ch:packages} is how.

\section{A note on \texttt{Undeclared}}

When you accept a method that references a name nothing defines yet, the
system does not refuse---it creates the binding, holding \ct{nil}, and
remembers it in \ct{Undeclared}. Later definition of that name fills the
binding in and everything works.

That is what makes mutually-referencing classes possible to type in any order.
It also means a typo in a global name compiles quietly and fails at run time
with \ct{nil doesNotUnderstand}. When a method behaves as though a global is
missing, look in \ct{Undeclared}:

\begin{st}
Undeclared
Undeclared keys
\end{st}

The bootstrap reports the same thing at build time---\ct|80 undeclared globals| is a line you will see on every build, and the lower-case ones in
that list are flagged separately because a lower-case undeclared global is
almost always a source bug.
